% !TEX encoding = UTF-8
% !TEX TS-program = pdflatex
% !TEX root = ../tesi.tex

\section{Tecnologie}\label{sec:tecnologie}
Di seguito viene data una panoramica delle tecnologie e degli strumenti utilizzati.

\subsection*{CFR - decompilatore java}
CFR è il decompilatore utilizzato per trasformare il codice bytecode \textit{.class} in codice sorgente \textit{.java};
\url{http://www.benf.org/other/cfr/index.html}

\subsection*{Apk Tool}
Apktool è un tool per effettuare il reverse engineering delle applicazioni Android.
Può decodificare le risorse contenute nell'APK in una forma quasi uguale a quelli originali, può anche di ricostruire l'APK dopo le modifiche alle risorse modificate.
Nel tool APAT è stato utilizzato per poter decompattare un'APK in questo modo permette di avere a disposizione i file che compongono l'APK originale.
L'apktool permette anche di ricomporre un'apk utilizzando le risorse ottenute precedentemente.
\url{https://ibotpeaches.github.io/Apktool/}
\begin{figure}[H]
    \centering
    \includegraphics[width=6cm, height=3cm]{./immagini/apktool.png}
    \caption{Icona Apktool}\label{fig:apktool}
\end{figure}

\subsection*{Dex2Jar}
Questo strumento permette di trasformare i file \textit{.dex} in formato \textit{.jar} in modo che possa essere letta attraverso un software di compressione, come 7Zip, per poter accedere quindi ai file di tipo \textit{.class}.
\url{https://github.com/pxb1988/dex2jar}

\subsection*{Java}
In informatica, Java\footcite{womak:effective-java} è un linguaggio di programmazione ad alto livello, orientato agli oggetti e tipizzazione statica, che si appoggia sull'omonima piattaforma software di esecuzione, specificamente progettato per essere il più possibile indipendente dalla piattaforma hardware di esecuzione.
Nell'APAT è stato utilizzato per la stesura del codice, in particolare, per l'interfaccia grafica è stata utilizzata la libreria di java    \gls{swing}.
\begin{figure}[H]
    \centering
    \includegraphics[width=8cm, height=5cm]{./immagini/java.jpg}
    \caption{Icona Java}\label{fig:java}
\end{figure}

\subsection*{Json}
JSON (JavaScript Object Notation) è un semplice formato per lo scambio di dati.
Per le persone è facile da leggere e scrivere, mentre per le macchine risulta facile da generare e analizzarne la sintassi.
JSON è basato su due strutture:
\begin{itemize}
    \item un insieme di coppie nome/valore, in diversi linguaggi, questo è realizzato come un oggetto, un record, uno struct, un dizionario, una tabella hash, un elenco di chiavi o un array associativo.
    \item un elenco ordinato di valori, nella maggior parte dei linguaggi questo si realizza con un array, un vettore o una sequenza.
\end{itemize}
Nell'APAT è stato utilizzato per contenere i dati di stato del tool quando viene chiuso e alcuni dati che indicano i percorsi ai tool di terze parti.
\url{https://www.json.org/json-it.html}
\begin{figure}[H]
    \centering
    \includegraphics[width=5cm, height=5cm]{./immagini/json.png}
    \caption{Icona JSON}\label{fig:json}
\end{figure}

\subsection*{XPath}\label{subsec:xpath}
In informatica XPath è un linguaggio, parte della famiglia \gls{xml}, che permette d'individuare i nodi all'interno di un documento XML. Le espressioni XPath, a differenza delle espressioni XML, non servono a identificare la struttura di un documento, bensì a localizzarne con precisione i nodi.
Nell'APAT è stato utilizzato per poter estrapolare alcuni dati necessari per poter effettuare il dump dei dati dall'area di storage dell'app e ottenere le informazioni per poter decompilare solamente presenti nel package dell'applicazione.


\section{Strumenti}\label{sec:strumenti}

\subsection*{Intellij IDEA}
L'\gls{ide} utilizzato per la scrittura del codice in Java sviluppato da JetBrains.
\url{https://www.jetbrains.com/idea/}
\begin{figure}[H]
    \centering
    \includegraphics[width=5cm, height=5cm]{./immagini/intellij.png}
    \caption{Icona Intellij Idea}\label{fig:intellij}
\end{figure}

\subsection*{Android Emulator}
Lo strumento che permette di avviare degli emulatori Android, conosciuti anche come \gls{avd}.
È stato utilizzato per eseguire l'applicazione ricompilata, per poter fare il dump dei dati presenti nell'area di storage dell'applicazione, come i file JSON, XML e SQLite, o come il file delle attività di rete.
\url{https://developer.android.com/studio/run/emulator}
\begin{figure}[H]
    \centering
    \includegraphics[width=4cm, height=4cm]{./immagini/emulator.png}
    \caption{Icona Android Emulator}\label{fig:emulator}
\end{figure}

\subsection*{Maven}
Maven è un tool di \gls{buildautomation}, basato sul concetto di Project Object Model (POM), maven può gestire i build, fare il report e la documentazione da un singolo centro di informazione.
\url{https://maven.apache.org/}
\begin{figure}[H]
    \centering
    \includegraphics[width=6cm, height=3cm]{./immagini/maven.png}
    \caption{Icona Maven}\label{fig:maven}
\end{figure}
Nel caso dell'APAT è stato utilizzato per la build del tool, generando in questo modo il file jar.

\subsection*{Astah UML}
Astah è uno strumento di modellazione \gls{uml}. Permette di creare vari diagrammi UML, come quelli di classe, di package, di sequenza e di attività.
\url{https://astah.net/products/astah-community/}
\begin{figure}[H]
    \centering
    \includegraphics[width=6cm, height=3cm]{./immagini/astah.png}
    \caption{Icona Astah UML}\label{fig:astah}
\end{figure}

\subsection*{Git - GitLab}
Gitlab è una piattaforma web open-source che permette la gestione di repository e di funzioni d'issue tracking system.
Questa è un'istanza privata dell'azienda Imola Informatica S.p.A..
\url{https://git.imolinfo.it/users/sign_in}
\begin{figure}[H]
    \centering
    \includegraphics[width=6cm, height=3cm]{./immagini/gitlab.png}
    \caption{Icona GitLab}\label{fig:gitlab}
\end{figure}
